% !TeX root = ../tesis.tex

% -----------------------------------------------------------------------
\section{Marco Teórico: Dispersión, Fragmentación y Equidad}
\label{sec:marco-teorico-dispersion}
% -----------------------------------------------------------------------

La consecuencia directa del modelo vial radial fue la expansión desordenada de
la
mancha urbana en un proceso denominado \termino{dispersión}. Sin embargo, para
los fines de esta tesis, el concepto más relevante es la \termino{fragmentación
urbana}, ya que denota una ruptura funcional y social que va más allá de la
simple
extensión física del territorio.

\citet{padilla2007} argumenta que la fragmentación se manifiesta en una
dislocación
de la cohesión urbana y social, donde las nuevas áreas periféricas carecen de
una
adecuada conexión funcional con el centro y entre sí. Esta dislocación se
traduce en
una desigualdad de acceso que perpetúa las brechas socioespaciales. Por lo
tanto,
el problema a resolver en el Anillo Periférico es la fragmentación territorial
resultante
del modelo radial, que impide la movilidad equitativa para las alcaldías del
sur.

% -----------------------------------------------------------------------
\subsection{Limitaciones del Modelo Radial}
\label{subsec:limitaciones-radial}
% -----------------------------------------------------------------------

La planificación del transporte público en la \cdmx{} siguió el esquema descrito
por
\citet{molinero2002}, donde las rutas radiales se priorizan para conectar
áreas
periféricas con un Distrito Central de Negocios (\textsc{cbd}, por sus siglas en
inglés). Este esquema predominó durante la segunda mitad del siglo XX, cuando la
ciudad experimentó un crecimiento urbano desmedido. Sin embargo, en ciudades
con más de 300,000 habitantes este tipo de rutas comienza a ser ineficiente, ya
que
concentra los movimientos en el centro y no considera las necesidades de
desplazamiento entre otras áreas urbanas \citep{molinero2002}.

% -----------------------------------------------------------------------
\subsection{Transiciones de Fase en Redes Urbanas}
\label{subsec:transiciones-fase}
% -----------------------------------------------------------------------

A la fecha de realización de este estudio (2025), no existe en la \cdmx{} alguna
ruta
circular que sirva de soporte para el transporte público masivo colectivo. A
pesar de
que la Ciudad cuenta con una vialidad de trazo circular (el Anillo Periférico),
esta
se ha orientado históricamente al automóvil privado. Esta ausencia representa
una
\termino{brecha topológica} en la red: la inexistencia de conexiones
transversales
impide que el sistema alcance una mayor eficiencia ante el crecimiento continuo
de
la demanda.

% -----------------------------------------------------------------------
\subsection{Minimización de Transferencias en Redes Anillo-Radiales}
\label{subsec:minimizacion-transferencias}
% -----------------------------------------------------------------------

Desde la perspectiva de la teoría de redes, \citet{owais2020} demuestra que
la
incorporación de una geometría anillar a un sistema predominantemente radial
reduce
el número de transbordos necesarios para conectar zonas periféricas entre sí.
Esta
minimización de transferencias no solo mejora los tiempos de traslado, sino que
incrementa la percepción de calidad del servicio por parte del usuario, al
reducir la
fricción del viaje multimodal que caracteriza actualmente la movilidad en las
alcaldías
del sur de la \cdmx{}.

